%----------------------------------------------------------------------------------------------------
% START OF CHAPTER
%----------------------------------------------------------------------------------------------------
\part{More Grammar} \label{M5-MoreGrammar}

\chapter{Modal Verbs: \itb{potere, volere, dovere}}

Modal verb conjugates as normal but the verb that follows stays in the \textbf{infinitive}.

\section{\itb{Potere}: can / to be able to}
\iti{posso, puoi, può, possiamo, potete, possono}

\section{\itb{Volere}: to want}
\iti{voglio, vuoi, vuole, vogliamo, volete, vogliono}

\textbf{Warning!} \iti{Vorrei} (I would like) is the \textbf{polite} conditional form; it is much more polite than \iti{voglio} for requests.

\section{\itb{Dovere}: must / to have to}
\iti{devo, devi, deve, dobbiamo, dovete, devono}

\section{\itb{Sapere} vs. \itb{potere}}
\begin{itemize}[nosep]
    \item \iti{Posso nuotare?} - Can I swim? (permission/physical ability)
    \item \iti{So nuotare.} - I can swim. (I know how to)
\end{itemize}

\section{Negation with Modals}
\iti{Non} goes before the modal: \iti{Non posso venire. Non voglio studiare. Non devo lavorare.}

\section{Pronouns with Modals}
Pronouns attach to the \itb{infinitive}. (They can, grammatically, go at the start of the phrase too, but the former is considered more natural):
\begin{itemize}[nosep]
    \item \iti{Posso \itb{aiutarti}} | \iti{Ti posso aiutare.} (both correct)
    \item \iti{Può \itb{dirmi}...?}  | \iti{Mi può dire...?}
\end{itemize}

\chapter{Reflexive Verbs}
Reflexive verbs describe actions done \itb{to oneself}. Their infinitives end in \itb{-si}.

Common reflexives: \iti{alzarsi, vestirsi, lavarsi, divertirsi, chiamarsi, sedersi}

\section{Reflexive Pronouns}
\begin{table}[h]
    \centering
    \begin{tabular}{r l}
        Person & Pronoun \\
        \hline
        io      & \itb{mi} \\
        tu      & \itb{ti} \\
        lui/lei & \itb{si} \\
        noi     & \itb{ci} \\
        voi     & \itb{vi} \\
        loro    & \itb{si} \\
    \end{tabular}
\end{table}

\section{Present Tense: pronoun before verb}
\iti{mi alzo, ti alzi, si alza, ci alziamo, vi alzate, si alzano}

\section{With Modals}
Pronoun attaches to infinitive (drop -e, keep -r):
\begin{itemize}[nosep]
    \item \iti{alzare → alzar + mi = \itb{alzarmi}}
    \item \iti{Voglio alzarmi presto.}
    \item \iti{Devi vestirti adesso!}
\end{itemize}


\section{Perfect Tense: reflexives always use \itb{essere}}
\textbf{reflexive pronoun + \iti{essere} + past participle (agrees with subject)}

\begin{table}[h]
    \centering
    \begin{tabular}{r l}
        Person & Example (alzarsi) \\
        \hline
        io (m)         & \itb{mi sono alzato}  \\
        io (f)         & \itb{mi sono alzata}  \\
        lui            & \itb{si è alzato}     \\
        lei            & \itb{si è alzata}     \\
        noi (m/mixed)  & \itb{ci siamo alzati} \\
        noi (f)        & \itb{ci siamo alzate} \\
        loro (m/mixed) & \itb{si sono alzati}  \\
        loro (f)       & \itb{si sono alzate}  \\
    \end{tabular}
\end{table}

\textbf{Note} Never drop the reflexive pronoun: \itb{si sono alzate}, \textbf{not} \textit{sono alzate}.

\chapter{The Perfect Tense with \itb{avere}: Passato Prossimo}

Used for completed past actions. Structure: \itb{avere} \textit{[conjugated]} + past participle.

\section{Forming the Past Participle}
\begin{table}[h]
    \centering
    \begin{tabular}{r | c l }
        Verb family & Ending & Example        \\
        \hline
        \itb{-are}  & \itb{-ato}  & \iti{parlare → parlato} \\
        \itb{-ere}  & \itb{-uto}  & \iti{vedere → veduto} \\
        \itb{-ire}  & \itb{-ito}  & \iti{dormire → dormito} \\
    \end{tabular}
\end{table}

\section{Common Irregular Past Participles}
\begin{table}[h]
    \centering
    \begin{tabular}{ r l | r l }
        Infinitive & Past participle & Infinitive & Past participle  \\
        \hline
        \iti{vedere} & \iti{visto} & \iti{fare} & \iti{fatto} \\
        \iti{leggere} & \iti{letto} & \iti{dire} & \iti{detto} \\
        \iti{scrivere} & \iti{scritto} & \iti{bere} & \iti{bevuto} \\
        \iti{prendere} & \iti{preso} & \iti{aprire} & \iti{aperto} \\
        \iti{perdere} & \iti{perso} & \iti{mettere} & \iti{messo} \\
        \iti{chiedere} & \iti{chiesto} & \iti{rispondere} & \iti{risposto} \\
        \iti{venire} & \iti{venuto} & \iti{uscire} & \iti{uscito} \\
    \end{tabular}
\end{table}

\section{Past Participle Agreement with Pronouns}
When a direct object pronoun (\iti{lo, la, li, le}) precedes \iti{avere}, the past participle agrees:
\begin{itemize}[nosep]
    \item \iti{le chiavi} → \iti{\itb{le} ho viste} (\iti{viste} feminine plural)
    \item \iti{il libro} → \iti{\itb{l'}ho letto} (\iti{letto} masculine singular)
\end{itemize}




    \chapter{The Perfect Tense with \iti{Essere}}
    Verbs of \textbf{motion and change of state} use \iti{essere}. The past participle \textbf{agrees with the subject}.

    \section{Key Verbs Using \iti{Essere}}
    \begin{table}[h]
        \centering
        \begin{tabular}{c | c l}
            Infinitive     & Participle     & Meaning    \\
            \hline
            \iti{andare}   & \iti{andato}   & to go      \\
            \iti{venire}   & \iti{venuto}   & to come    \\
            \iti{entrare}  & \iti{entrato}  & to enter   \\
            \iti{uscire}   & \iti{uscito}   & to exit    \\
            \iti{arrivare} & \iti{arrivato} & to arrive  \\
            \iti{partire}  & \iti{partito}  & to depart  \\
            \iti{tornare}  & \iti{tornato}  & to return  \\
            \iti{stare}    & \iti{stato}    & to stay    \\
            \iti{essere}   & \iti{stato}    & to be      \\
            \iti{rimanere} & \iti{rimasto}  & to remain  \\
            \iti{nascere}  & \iti{nato}     & to be born \\
            \iti{morire}   & \iti{morto}    & to die     \\
        \end{tabular}
    \end{table}

    All \textbf{reflexive verbs} also use \iti{essere}.

    \section{Past Participle Agreement}
    \begin{table}[h]
        \centering
        \begin{tabular}{c c l}
            Subject        & Ending     & Example    \\
            \hline
            Masc. sing.    & \iti{-o}   & \iti{Mario è andato} \\
            Fem. sing.     & \iti{-a}   & \iti{Maria è andata} \\
            Masc/mixed pl. & \iti{-i}   & \iti{Mario e Luigi sono andati} \\
            Fem. pl.       & \iti{-e}   & \iti{Le ragazze sono andate} \\
        \end{tabular}
    \end{table}

\section{Recap: \iti{Avere} or \iti{Essere}?}
\begin{itemize}[nosep]
    \item Verb of motion/change + \textbf{no direct object} → \iti{essere}
    \item Verb with a \textbf{direct object} → \iti{avere}
\end{itemize}


\chapter{Negation, Questions and Sentence Structure}

\section{Basic Negation}
\itb{Non} goes before the verb: \iti{Non parlo italiano. Non ho mangiato. Non mi sono alzato.}
In the perfect tense, put \iti{non} before \iti{avere} or \iti{essere}. \iti{Non ho mangiato. Non abbiamo fatto niente. Non sono andate a Roma.}

\section{Negative Word Pairs}
\begin{table}[h]
    \centering
    \begin{tabular}{c c l}
        Positive        & Negative     & Meaning    \\
        \hline
        \iti{qualcuno} & \iti{nessuno} & someone → nobody            \\
        \iti{qualcosa} & \iti{niente / nulla} & something → nothing  \\
        \iti{sempre} & \iti{non... mai} & always → never             \\
        \iti{già} & \iti{non... ancora} & already → not yet          \\
        \iti{ancora} & \iti{non... più} & still → no longer          \\
        \iti{anche} & \iti{non... neanche/nemmeno} & also → not even \\
    \end{tabular}
\end{table}

\section{Placement Rules}

\textbf{Simple tenses}: \iti{non} + verb + negative word
\begin{itemize}[nosep]
    \item \iti{Non vado \itb{mai} al cinema.}
    \item \iti{Non lavoro \itb{più} lì.}
\end{itemize}

\textbf{Perfect tense}: \iti{non} + auxiliary + negative word + participle
\begin{itemize}[nosep]
    \item \iti{Non sono \itb{mai} uscito.}
    \item \iti{Non ho \itb{ancora} finito.}
\end{itemize}

\section{Word Order: negative words in perfect tense}
Use \iti{non} + auxiliary + negative word + past participle:
\begin{itemize}[nosep]
    \item \iti{non ho \itb{mai} mangiato}
    \item \iti{non ho \itb{ancora} finito}
    \item \iti{non ho \itb{neanche} studiato}
\end{itemize}

\textbf{Negative word at start of sentence}: drop the \iti{non}:
\begin{itemize}[nosep]
    \item \iti{\itb{Nessuno} è venuto.}
    \item \iti{\itb{Niente} da fare!}
\end{itemize}

\section{\iti{Nessuno} as Subject \textit{vs.} Object}
\begin{itemize}[nosep]
    \item Subject (start): \iti{Nessuno mi ha chiamato.} (no \iti{non})
    \item Object (after verb): \iti{Non ho visto nessuno.} (needs \iti{non})
\end{itemize}

\section{\iti{Niente/qualcosa} + \iti{di} + adjective}
Always insert \itb{di} between \iti{niente/qualcosa} and an adjective:
\begin{itemize}[nosep]
    \item \iti{niente \itb{di} interessante} | \iti{qualcosa \itb{di} bello}
\end{itemize}

\section{``Have you ever'' = \iti{non... mai}}
\iti{Non sei mai stato in Italia?}: Have you ever been to Italy?

\chapter{Question Words}
\begin{table}[h]
    \centering
    \begin{tabular}{c c l}
        Word      & Meaning   & Example             \\
        \hline
        \iti{che / cosa / che cosa}  & what          & \iti{Cosa fai?}         \\
        \iti{chi}                    & who           & \iti{Chi viene?}        \\
        \iti{dove}                   & where         & \iti{Dove sei andato?}  \\
        \iti{quando}                 & when          & \iti{Quando parti?}     \\
        \iti{come}                   & how           & \iti{Come stai?}        \\
        \iti{perché}                 & why / because & \iti{Perché non vieni?} \\
        \iti{quanto/a}               & how much      & \iti{Quanto costa?}     \\
        \iti{quanti/e}               & how many      & \iti{Quanti anni hai?}  \\
        \iti{quale/quali}            & which         & \iti{Quale preferisci?} \\
    \end{tabular}
\end{table}

\section{Question Word Order}
The form is \textbf{Question word} + \textbf{verb} + \textbf{subject}:

\begin{itemize}[nosep]
    \item \iti{Dove sono andati i tuoi genitori?}
    \item \textit{Dove i tuoi genitori sono andati?} \textbf{Incorrect!}
\end{itemize}

\section{\iti{Chi} vs \itb{Che}}
\begin{itemize}[nosep]
    \item \itb{Chi} = who in questions / whoever (no noun before it)
    \item \itb{Che} = who/which/that in relative clauses (noun before it)
\end{itemize}

Examples: \iti{itb{Chi} parla italiano?} - Who speaks Italian? (question). Versus
\iti{Gli studenti \itb{che} parlano italiano...} - The students who speak Italian... (relative clause)

%----------------------------------------------------------------------------------------------------
% END OF CHAPTER
%----------------------------------------------------------------------------------------------------